\chapter{Introduction}
\label{ch:introduction}

\section{Motivation}
\label{sec:motivation}

Almost every major securities and derivatives exchange in the world
organises trading as a continuous double auction (CDA), implemented as a
continuous limit order book (CLOB). Orders arrive over a network, are
placed in a queue, and are processed one at a time in the order in which
the exchange received them. Whenever an incoming order crosses the
opposite side of the book, a trade is executed immediately at the
resting limit price. The rule is simple, it is easy to explain to
participants, and it has the appealing property that a trader who is
willing to pay the current best price never has to wait.

The same rule has a structural consequence that is far less benign.
Because execution is serial and instantaneous, the economic value of a
trading opportunity accrues entirely to whoever reaches the matching
engine first. When public information changes the fair value of an
asset, the stale quotes resting in the book become mispriced for a short
interval, and the first participant to act either takes the stale quote
or cancels it. Nothing about the informational content of the event
determines who wins; only relative speed does. This turns a mechanism
that was designed to allocate goods according to willingness to pay into
one that partly allocates them according to willingness to invest in
latency reduction \citep{budish2015hft}.

The resulting competition is not costless. Microwave and hollow-core
fibre links, co-location, custom hardware, and kernel-bypass networking
absorb real resources without improving the informational efficiency of
prices in any first-order way. More importantly for the participants who
are not in the race, the risk of being picked off is priced into
quotes. A market maker who knows that a fraction of its fills will occur
precisely at the moment its quote has become stale must widen its spread
to break even on the remainder. 
\citet{aquilina2022quantifying} exploit
message-level data from the London Stock Exchange to measure this cost
directly. They identify the races themselves in the raw message stream,
show that they are extremely frequent and extremely short, and quantify
the transfer as an implicit tax on liquidity demanders. Their central
estimate implies that a meaningful share of the effective spread on a
liquid instrument is compensation for latency-driven adverse selection
rather than for inventory risk or order processing cost.

The frequent batch auction (FBA) is the market design response proposed
by \citet{budish2015hft}. Instead of processing orders one by one as they
arrive, the engine collects all orders that arrive within a short
discrete interval $\tau$, and at the end of the interval clears them
simultaneously against each other at a single uniform price. Two things
change at once. First, arrival order inside the interval no longer
matters: an order submitted one microsecond after the interval opens and
an order submitted one microsecond before it closes are treated
identically. Second, all trades in the interval settle at the same
price, so a participant cannot capture a better price simply by being
earlier in the queue. Speed retains value only across intervals, and the
marginal return to shaving another nanosecond off a network path
collapses to approximately zero. Competition is redirected from latency
to price.

The idea is not exotic. Almost every equity exchange already runs
uniform-price call auctions at the open and the close, and several
markets have used them intraday. What is new in the FBA proposal is the
frequency: an auction every few hundred milliseconds, or faster, so that
the delay imposed on ordinary traders is small enough to be irrelevant
while still being long enough to swamp the speed differences that fund
the arms race. In decentralised finance the argument arrives from a
different direction but converges on the same mechanism. Blockchain
settlement is already discrete, transaction ordering inside a block is
economically contestable, and batch clearing with a uniform price
removes the ordering rent that would otherwise be extracted by block
producers \citep{gong2023fbadex,walther2021multitoken}. Several deployed
systems, among them CoW Protocol, Loopring, and SPEEDEX, are
uniform-price batch auctions in all but name
\citep{dutta2025cow,wang2018loopring,ramsay2023speedex}.

\section{The Design Problem behind the Idea}
\label{sec:designproblem}

Stating that orders should be batched and cleared at a uniform price
does not yet specify a matching engine. It specifies a family of
matching engines. An exchange operator who sets out to implement one
immediately faces a series of decisions that the theoretical proposal
leaves open, and each of these decisions changes the strategic problem
that participants face.

How long should the interval be, and should its end be predictable? A
fixed 100~millisecond interval is simple to reason about and easy to
integrate with existing systems, but its boundary is a schedule, and a
schedule can be attacked: a sufficiently fast participant can wait until
the last possible moment, observe as much information as possible, and
submit or cancel just before the gate closes. Randomising the closing
time removes the target but complicates order management for everyone
else.

What should participants see while the auction is collecting orders? An
open book during the accumulation phase helps the market converge on a
sensible clearing price and gives participants the information they need
to compete, but it also broadcasts intent and invites the interval
equivalent of front-running. A sealed-bid design eliminates that
exposure entirely, at the price of clearing an auction in which nobody
knows what anybody else has done, which classical auction theory
predicts will lead to more conservative bidding
\citep{vickrey1961counterspeculation,klemperer1999auction}.

Should all order flow meet in a single pool? Retail flow is on average
uninformed, institutional flow on average is not, and a market maker who
cannot distinguish them must quote a price that reflects the mixture
\citep{glosten1985bid}. Segregating flow allows tighter quotes against
the benign part of it. Jump Trading's Dual Flow Batch Auction takes this
further by partitioning each batch into maker and taker intent groups
and running two independent auctions per interval, so that liquidity
providers never trade against each other.

What happens when the two sides do not match exactly at the clearing
price? Uniform-price clearing determines the price and the total
quantity, but not who among the participants on the long side of the
market gets filled. The rationing rule that resolves this is not a
technicality. Pro-rata allocation invites participants to inflate their
order sizes; time priority inside the batch reintroduces a micro-scale
speed race through the back door; random allocation is strategically
clean but unpredictable for participants who need certainty of
execution. \citet{haynes2020precedence} document how sensitive
participant behaviour is to precedence rules in existing derivatives
markets.

None of these questions has a single correct answer. Each represents a
trade-off, and the purpose of the theoretical part of this thesis is to
make those trade-offs explicit and to derive, from auction theory and
market microstructure theory, which behaviour each option should be
expected to induce.

\section{Problem Statement and Research Gap}
\label{sec:gap}

Two gaps motivate this work, one conceptual and one empirical.

\paragraph{The design space is discussed piecemeal.}
The literature on frequent batch auctions is largely a literature about
whether batching is a good idea, not about how a batching engine should
be built. \citet{budish2015hft} establish the equilibrium argument,
\citet{aquilina2022quantifying} quantify what the alternative costs,
\citet{zoican2021speed} and \citet{eibelshauser2022fba} study how
information and speed interact inside auctions, and
\citet{aldrich2020experiments} compare the two institutions in the
laboratory. Individual design dimensions appear in isolation: rationing
rules in the derivatives microstructure literature, disclosure in the
transparency literature \citep{madhavan1996transparency}, and clearing
formulations in the decentralised exchange literature
\citep{walther2021multitoken,ramsay2023speedex}. What is missing is a
consolidated treatment that decomposes the engine into its architectural
blocks, sets out the options for each, and traces each option to the
strategic behaviour it induces. Two further points inside this space are
particularly underdeveloped. First, in continuous markets the
maker--taker distinction is defined by arrival timing relative to book
state, but in a batch that clears simultaneously at one price that
definition dissolves, which means fee and rebate schemes must be
redefined rather than transplanted. Second, uniform-price clearing
mechanically produces price improvement relative to the submitted limit
price for every executed order except the marginal one, and the size and
distribution of that improvement depend directly on the clearing and
rationing rules chosen.

\paragraph{Field comparisons cannot hold order flow constant.}
The empirical literature comparing batch and continuous trading is
forced to compare different venues, different assets, or different time
periods. \citet{lee2026taiwan} exploit the Taiwan Stock Exchange's
transition from intraday call auctions to continuous trading, which is
about as clean as a natural experiment in this area can be, and yet the
participant population, the order submission strategies, and the
information environment all adjust together with the mechanism. This is
not a shortcoming of the studies; it is a property of the object. An
order stream exists only under the mechanism that produced it, so a
strictly like-for-like comparison in which one and the same stream is
executed once continuously and once in batches cannot be obtained from
field data at all. Simulation with two interchangeable engines is the
only design that holds the flow fixed and varies the mechanism.

\section{Research Questions}
\label{sec:rqs}

The thesis is organised around two research questions, one for each
part.

\begin{quote}
\textbf{RQ1 (theoretical).} How is a frequent batch auction engine
designed, and which design options exist for each of its architectural
components?
\end{quote}

RQ1 is answered by decomposing the engine into five design blocks ---
the batch interval, the information disclosure regime, order flow
segregation, the matching and rationing rule, and settlement and
residual handling --- by enumerating the options available within each
block, and by deriving from theory the strategic behaviour that each
option induces. The output is a structured description of the design
space together with a trade-off matrix.

\begin{quote}
\textbf{RQ2 (empirical).} How do a continuous double auction and a
frequent batch auction differ when both engines are fed with the same
order flow?
\end{quote}

RQ2 is decomposed into three measurable dimensions:

\begin{itemize}
  \item \textbf{RQ2.1 Liquidity and transaction costs.} Which mechanism
  lets participants trade at a lower cost and with more depth
  available?
  \item \textbf{RQ2.2 Price discovery and market quality.} Which
  mechanism produces a more stable, less dispersed price series?
  \item \textbf{RQ2.3 Execution, allocation and order-arrival timing.}
  Who gets filled, at what cost, and how does the batch mechanism
  reshape execution outcomes for the same unchanged order flow?
\end{itemize}

Each dimension is operationalised by a set of metrics defined in
\cref{sec:metrics}, and every metric is tagged with the minimum order
book data level it requires.

\section{Contribution}
\label{sec:contribution}

The thesis makes three contributions. It consolidates the FBA design
space into five blocks and links each available option to a theoretical
prediction about induced behaviour, which turns a scattered discussion
into a checklist an exchange operator can work through. It implements
two matching engines behind a single shared interface, so that the
mechanism is the only treatment variable and every internal quantity ---
including submitted limit prices of orders that never execute, the exact
rationing outcome, and the residual left unmatched after each batch ---
is directly observable rather than inferred. And it applies this
apparatus to a recently released Level~4 order book dataset from the
Hyperliquid exchange \citep{albers2026openbook}, which carries
participant attribution at nanosecond precision and therefore supports
per-participant allocation metrics that conventional Level~2 or Level~3
feeds cannot.


